\section{Method}
\label{sec:method}
\begin{figure*}[t]
\centering
\begin{tikzpicture}[x=1cm,y=1cm,
box/.style={draw=black!45,rounded corners=2pt,align=center,minimum height=1.12cm,font=\small,inner sep=6pt},
flow/.style={-{Latex[length=2mm]},thick,draw=black!65},
sup/.style={flow,dashed,draw=orange!75!black}]
\node[box,fill=blue!7,text width=2.5cm] (build) at (0,0) {Build LiDAR\\endpoints + ray history};
\node[box,fill=blue!7,text width=2.6cm] (shape) at (3.8,0) {Canonical surface\\fixed anchors + 4 children};
\node[box,fill=green!8,text width=2.7cm] (mass) at (7.8,0) {Frozen geometry\\learned F / O / U evidence};
\node[box,fill=green!8,text width=2.4cm] (return) at (11.6,0) {Categorical measure\\interpolated ray return};
\draw[flow] (build)--node[above,font=\scriptsize]{set decoder}(shape);
\draw[flow] (shape)--node[above,font=\scriptsize]{evidence heads}(mass);
\draw[flow] (mass)--node[above,font=\scriptsize]{normalize}(return);
\node[box,fill=orange!8,text width=5.3cm,minimum height=.83cm] (target) at (3.8,-1.72)
{Disjoint target sweeps: set / frame / ray supervision};
\draw[sup] (target.north)--(shape.south);
\draw[sup] (target.north east)--++(0,.25)-|(mass.south);
\node[box,fill=purple!7,text width=3.0cm,minimum height=.83cm] (pose) at (9.1,-1.72)
{Read-only SE(3) pose\\canonical $\leftrightarrow$ world};
\draw[flow] (pose.east)-|(return.south);
\node[font=\scriptsize,align=center,text=black!70] at (4.5,-2.65)
{Physical state $P$ \quad $\xrightarrow{\;\mathrm{stop\mbox{-}gradient}\;}$ \quad visual state $V$ \quad $\xleftarrow{\;\mathrm{image\ loss}\;}$ \quad RGB};
\node[font=\scriptsize,align=center,text=black!70] at (10.8,-2.65)
{Separate static Background};
\end{tikzpicture}
\caption{Overview of Evidential Actor Surfaces. Motion-compensated build observations produce a canonical surface.
Disjoint targets supervise geometry and continuous free/occupied/unknown (F/O/U) evidence in stages.
Frozen evidence weights a categorical ray measure. Rigid pose and appearance have separate state ownership.
Dashed arrows denote supervision.}
\label{fig:eas-overview}
\end{figure*}


\subsection{Motion-Compensated Surface Supervision}

Annotated sensor, ego, and Actor poses map each endpoint into a canonical frame,
\begin{equation}
x_{it}=T_{it}^{-1}T_{\mathrm{ego},t}T_{\mathrm{sensor}}x_t^{\mathrm{sensor}}.
\end{equation}
We transform its sensor origin by the same chain. Build sweeps produce fixed observed anchors $\mathcal A_i$
and completion seeds; separate sweeps produce the target set $\mathcal Y_i$.
Observed anchors comprise retained endpoints and endpoints projected onto build-supported geometry.

A pretrained relocation model supplies parent center $\bar x_j$. A set decoder combines build features with
four learned slot embeddings and predicts bounded residuals and isotropic scales:
\begin{equation}
c_{jk}=\bar x_j+\Delta x_{jk},\quad
S_i=\mathcal A_i\cup\{(c_{jk},s_{jk}):k=1,\ldots,4\}.
\label{eq:child-surface}
\end{equation}
Residuals are bounded by $0.25$\,m per coordinate and scales lie in $[0.02,0.40]$\,m.
We write $\mathcal X_i$ for the centers of $S_i$.
The decoder uses Actor dimensions and canonical geometric evidence. Its target is a set, allowing several
children to cover distinct parts of a local surface.

\paragraph{Geometry and frame coverage.}
Let $\widetilde{\mathcal L}$ denote a loss normalized by its detached parent-reference value.
We combine symmetric nearest-neighbor distance, local point-to-plane distance, and log-scale regression:
\begin{equation}
\mathcal L_G=\widetilde{\mathcal L}_{\mathrm{set}}
+.25\widetilde{\mathcal L}_{\mathrm{plane}}
+.10\widetilde{\mathcal L}_{\mathrm{scale}}
+\widetilde{\mathcal L}_{\mathrm{frame}}.
\label{eq:physical-loss}
\end{equation}
Local planes and scale targets come from eight target neighbors. The frame term averages directed
target-to-surface distance equally over target frames,
\begin{equation}
\mathcal L_{\mathrm{frame}}=
\frac{1}{|\mathcal F_i|}\sum_{t\in\mathcal F_i}
\frac{1}{|\mathcal Y_{it}|}\sum_{y\in\mathcal Y_{it}}\min_{x\in\mathcal X_i}\|y-x\|_2.
\label{eq:frame-loss}
\end{equation}
This prevents densely sampled frames from dominating canonical coverage.

\paragraph{Differentiable ray supervision.}
We sample 48 depths along each training ray and alpha-compose Gaussian support into an expected depth
$\bar d=\sum_k T_k\alpha_k d_k+T_{\mathrm{end}}d_{\mathrm{far}}$.
Here $d_{\mathrm{far}}=d^\star+0.50$\,m supplies the residual no-hit depth.
The ray objective is
\begin{equation}
\mathcal L_R=
\widetilde{\operatorname{Huber}}(\bar d,d^\star)
+.5\,\widetilde{[d^\star-\tau-\bar d]_+}.
\label{eq:smooth-ray-loss}
\end{equation}
The second term penalizes a predicted depth inside the measured pre-hit interval. This is a smooth
depth surrogate; the literal beam-tube minimum is used for evaluation.
During physical fine-tuning, PCGrad~\cite{yu2020pcgrad} projects conflicting geometry and ray gradients.
All children are emitted at deployment, followed by a fixed $0.06$\,m voxel deduplication.

\subsection{Producer-Evidential Return Measure}

We next fix the geometry and learn how its primitives contribute to sensor termination.
For each anchor, the producer retains source-frame and ray provenance, projection displacement,
canonical hit count, temporal/view support, and continuous build evidence.
Children inherit parent evidence together with their local geometry.
Separate PointNet-style heads predict
$m_j=(m_{F,j},m_{O,j},m_{U,j})$ by a three-way softmax.
Target-ray votes supervise these masses with soft cross-entropy: observations before a hit contribute
free evidence, endpoints contribute occupied evidence, and unsupported or occluded regions contribute unknown evidence.

For a completion child, let $n_j$ count positive-depth ray opportunities inside the beam tube,
and let $n_{F,j}$ and $n_{O,j}$ count pre-hit and endpoint votes. Its target is
\begin{equation}
m_j^\star=\left(\frac{n_{F,j}}{n_j},\frac{n_{O,j}}{n_j},
1-\frac{n_{F,j}+n_{O,j}}{n_j}\right),\quad n_j>0,
\label{eq:child-votes}
\end{equation}
with $(0,0,1)$ when no ray supports the child. Both free and occupied coordinates can be positive,
preserving disagreement across observations. The evidence loss is
$\mathcal L_E=-\sum_j\sum_{a\in\{F,O,U\}}m_{a,j}^\star\log m_{a,j}$,
averaged over the supervised primitives.

\paragraph{Staged evidence learning.}
The anchor head is trained with evidential and ordered-transmittance return losses.
With that head frozen, the child head is trained with the same objectives and then fine-tuned with
a pre-hit survival term, the integrated child optical thickness before $d^\star-\tau$.
These stages learn evidence on fixed centers and scales.
We subsequently freeze both heads and change the return composition to a categorical measure.
Thus the final composition is a frozen representation experiment; it does not require joint categorical fine-tuning.

\paragraph{Categorical composition.}
For 64 ordered samples $d_k$ inside the Actor box, define
\begin{equation}
\kappa_{kj}=\exp\!\left(-\frac{\|o+d_kv-c_j\|_2^2}{2s_j^2}\right),
\qquad w_j=m_{O,j}.
\end{equation}
The joint depth--primitive distribution is
\begin{equation}
q_{kj}=\frac{w_j\kappa_{kj}}{\sum_{a,b}w_b\kappa_{ab}},
\qquad p_k=\sum_jq_{kj},\quad F_k=\sum_{\ell\leq k}p_\ell.
\label{eq:categorical-return}
\end{equation}
We interpolate within the first CDF bin crossing $1/2$ to obtain the return $\widehat d$.
The measure conditions on a return inside the annotated box.
Multiplying all primitive weights by the same positive constant leaves $p_k$ unchanged.
This removes a global optical-thickness degree of freedom; changing the relative mass of a primitive
family still changes the distribution.

\subsection{Component Responsibility and Attenuation}
\label{sec:attenuation}

For an evaluation boundary $b=d^\star-\tau$, interpolate $F(b)$ with the same bins as the median.
At interior boundaries with positive bin mass,
\begin{equation}
\widehat d<b\quad\Longleftrightarrow\quad F(b)>\tfrac12.
\label{eq:categorical-safety-boundary}
\end{equation}
Because the measure is additive over primitives,
$F(b)=F_{\mathcal A}(b)+F_{\mathrm{child}}(b)$.
This decomposes the deployed early event into observed-anchor and completion contributions.

Let $D_j=\sum_k\kappa_{kj}$, let $N_j(b)$ be the interpolated pre-boundary kernel mass, and define
$C_j=N_j/D_j$ and $r_j=w_jD_j/\sum_a w_aD_a$.
Then
\begin{equation}
\frac{\partial F(b)}{\partial\log w_j}=r_j(C_j-F(b)).
\label{eq:attenuation-sign-boundary}
\end{equation}
For a finite attenuation $w_j\mapsto vw_j$, $0<v<1$,
\begin{equation}
F_v(b)-F(b)=
\frac{(1-v)r_j\bigl(F(b)-C_j\bigr)}{1-(1-v)r_j}.
\label{eq:finite-attenuation-boundary}
\end{equation}
The denominator is positive. For a component of positive mass, attenuation decreases $F(b)$ exactly
when $C_j>F(b)$; it increases $F(b)$ when $C_j<F(b)$.
The same identity holds for uniform attenuation of a family after aggregating its mass.
Its sign is independent of attenuation magnitude for that fixed component or family.
Indeed, the attenuated measure is $F_v=(F-(1-v)r_jC_j)/(1-(1-v)r_j)$;
subtracting $F$ gives Eq.~\ref{eq:finite-attenuation-boundary}.
We use target-defined $b$ for analysis.

\subsection{Rigid Composition and Visual Ownership}

The canonical energy $e_i(x)=\log\sum_j w_j\exp(-\|x-c_j\|^2/(2s_j^2))$ is placed in the world by
\begin{equation}
Q_{it}(x)=e_i\!\left(R_{it}^{\top}(x-t_{it})\right),
\qquad T_{it}=(R_{it},t_{it}).
\label{eq:authority-query}
\end{equation}
For a common rigid transform $g$, inverse querying gives $Q_{gT}(gx)=Q_T(x)$.
Rigid pose changes location while preserving canonical pairwise distances.
Static Background has a separate scene owner.

The visual layer receives a detached physical carrier $\operatorname{sg}(P_i)$ and owns its render geometry,
spherical harmonics, and opacity. Image losses update only this sibling state.
Since the physical query reads $(P_i,T_{it})$, its derivative with respect to visual parameters is zero.
This interface lets appearance capacity grow independently of the physical surface.
